% GENERATED FILE. Edit canonical lessons, then run npm run generate:books.
% canonical-source-hash: fnv1a64:95caae92e199934b
% canonical-lessons: SA-C553-skhalanam, SA-C553-hanta, SA-C553-pariraksati, SA-C553-viparyayah, SA-C553-upalambhah

\chapter{If Only}
\label{ch:sa-if-only}
\hlchaptermodality{553}

\begin{chapteropening}
\textbf{By the end of this chapter:} \emph{I can say what would have happened if things had gone differently.}

Complete the last lesson of chapter 553: I can say what would have happened if things had gone differently.
\end{chapteropening}

\section[\texorpdfstring{\sk{स्खलनम्} (skhalanam)}{स्खलनम् (skhalanam)}]{\sk{स्खलनम्} (skhalanam) — a slip}

\label{lesson:SA-C553-skhalanam}



\begin{quote}
Before the new one: say the Sanskrit for makes use of, then the Sanskrit for helpful.
\end{quote}

\subsection*{You'll want to know: \sk{स्खलनम्}}
\textbf{\sk{स्खलनम्}} — \emph{skhalanam} — \textquotedblleft{}a slip\textquotedblright{}.

\begin{quote}
\sk{यदि} \sk{स्खलनं} \sk{न} \sk{अभविष्यत्}, \sk{तर्हि} \sk{सः} \sk{विजेता} \sk{अभविष्यत्।}
\end{quote}

\begin{quote}
\emph{yadi skhalanaṃ na abhaviṣyat, tarhi saḥ vijetā abhaviṣyat.}
\end{quote}

\begin{quote}
If there had been no slip, he would have been the victor.
\end{quote}

A neuter noun.

\begin{cousinweb}
\textbf{\sk{स्खलनम्}} (\emph{skhalanam}) comes from the root √skhal, \textquotedblleft{}to stumble, to trip\textquotedblright{}, with the ending \emph{-ana}; MW gives \emph{skhalana} as stumbling, tripping, faltering (MW).
\end{cousinweb}

\begin{grammarlens}[title={What you've built}]
One more word for saying what happens if.
\end{grammarlens}

\subsection*{Your turn}
\begin{itemize}
\raggedright
  \item \emph{Say it:} \emph{skhalanam}
  \item \emph{Say it:} \emph{skhalanam}, once more
  \item \emph{Say it:} \emph{yadi skhalanaṃ na abhaviṣyat, tarhi saḥ vijetā abhaviṣyat.} — If there had been no slip, he would have been the victor.
  \item \emph{Recall:} say the Sanskrit for makes use of, then the Sanskrit for helpful, then say \emph{skhalanam} again
\end{itemize}

\begin{tcolorbox}[breakable,colback=teal!4,colframe=teal!35!black,title={Before you move on}]
What does \sk{स्खलनम्} mean? (\textquotedblleft{}A slip\textquotedblright{}.) Say it once more.
\end{tcolorbox}
\section[\texorpdfstring{\sk{हन्त} (hanta)}{हन्त (hanta)}]{\sk{हन्त} (hanta) — come now; alas}

\label{lesson:SA-C553-hanta}



\begin{quote}
Before the new one: say the Sanskrit for helpful, then the Sanskrit for a slip.
\end{quote}

\subsection*{You'll want to know: \sk{हन्त}}
\textbf{\sk{हन्त}} — \emph{hanta} — \textquotedblleft{}come now; alas\textquotedblright{}.

\begin{quote}
\sk{हन्त}, \sk{यदि} \sk{अहं} \sk{ह्यः} \sk{अगमिष्यम्}, \sk{तर्हि} \sk{तम्} \sk{अद्रक्ष्यम्।}
\end{quote}

\begin{quote}
\emph{hanta, yadi ahaṃ hyaḥ agamiṣyam, tarhi tam adrakṣyam.}
\end{quote}

\begin{quote}
Alas, if I had gone yesterday, I would have seen him.
\end{quote}

\begin{cousinweb}
\textbf{\sk{हन्त}} (\emph{hanta}) is the particle MW gives for urging and for asking attention, \textquotedblleft{}come, see\textquotedblright{}, and in later use for grief, joy or pity; MW gives no derivation, so its origin is uncertain (MW).
\end{cousinweb}

\begin{grammarlens}[title={What you've built}]
One more word for saying what happens if.
\end{grammarlens}

\subsection*{Your turn}
\begin{itemize}
\raggedright
  \item \emph{Say it:} \emph{hanta}
  \item \emph{Say it:} \emph{hanta}, once more
  \item \emph{Say it:} \emph{hanta, yadi ahaṃ hyaḥ agamiṣyam, tarhi tam adrakṣyam.} — Alas, if I had gone yesterday, I would have seen him.
  \item \emph{Recall:} say the Sanskrit for helpful, then the Sanskrit for a slip, then say \emph{hanta} again
\end{itemize}

\begin{tcolorbox}[breakable,colback=teal!4,colframe=teal!35!black,title={Before you move on}]
What does \sk{हन्त} mean? (\textquotedblleft{}Come now; alas\textquotedblright{}.) Say it once more.
\end{tcolorbox}
\section[\texorpdfstring{\sk{परिरक्षति} (parirakṣati)}{परिरक्षति (parirakṣati)}]{\sk{परिरक्षति} (parirakṣati) — guards well; rescues}

\label{lesson:SA-C553-pariraksati}



\begin{quote}
Before the new one: say the Sanskrit for a slip, then the Sanskrit for come now.
\end{quote}

\subsection*{You'll want to know: \sk{परिरक्षति}}
\textbf{\sk{परिरक्षति}} — \emph{parirakṣati} — \textquotedblleft{}guards well; rescues\textquotedblright{}.

\begin{quote}
\sk{यदि} \sk{रक्षकः} \sk{न} \sk{आगमिष्यत्}, \sk{तर्हि} \sk{कः} \sk{बालं} \sk{पर्यरक्षिष्यत्}?
\end{quote}

\begin{quote}
\emph{yadi rakṣakaḥ na āgamiṣyat, tarhi kaḥ bālaṃ paryarakṣiṣyat?}
\end{quote}

\begin{quote}
If the guard had not come, who would have saved the child?
\end{quote}

Past (he …): \sk{सः} \sk{बालं} \sk{परिरक्षितवान्}, \emph{saḥ bālaṃ parirakṣitavān}.

Future: \sk{परिरक्षिष्यति}, \emph{parirakṣiṣyati}.

\begin{cousinweb}
\textbf{\sk{परिरक्षति}} (\emph{parirakṣati}) is built on \emph{pari}, \textquotedblleft{}all round\textquotedblright{}, and the root √rakṣ, \textquotedblleft{}to protect\textquotedblright{}, which also gives \emph{rakṣati}; MW gives \emph{parirakṣ} as to guard well, to rescue, to save (MW).
\end{cousinweb}

\begin{grammarlens}[title={What you've built}]
One more word for saying what happens if.
\end{grammarlens}

\subsection*{Your turn}
\begin{itemize}
\raggedright
  \item \emph{Say it:} \emph{parirakṣati}
  \item \emph{Say it:} \emph{parirakṣati}, once more
  \item \emph{Say it:} \emph{yadi rakṣakaḥ na āgamiṣyat, tarhi kaḥ bālaṃ paryarakṣiṣyat?} — If the guard had not come, who would have saved the child?
  \item \emph{Recall:} say the Sanskrit for a slip, then the Sanskrit for come now, then say \emph{parirakṣati} again
\end{itemize}

\begin{tcolorbox}[breakable,colback=teal!4,colframe=teal!35!black,title={Before you move on}]
What does \sk{परिरक्षति} mean? (\textquotedblleft{}Guards well; rescues\textquotedblright{}.) Say it once more.
\end{tcolorbox}
\section[\texorpdfstring{\sk{विपर्ययः} (viparyayaḥ)}{विपर्ययः (viparyayaḥ)}]{\sk{विपर्ययः} (viparyayaḥ) — a reversal; a mix-up}

\label{lesson:SA-C553-viparyayah}



\begin{quote}
Before the new one: say the Sanskrit for come now, then the Sanskrit for guards well.
\end{quote}

\subsection*{You'll want to know: \sk{विपर्ययः}}
\textbf{\sk{विपर्ययः}} — \emph{viparyayaḥ} — \textquotedblleft{}a reversal; a mix-up\textquotedblright{}.

\begin{quote}
\sk{यदि} \sk{दिशायाः} \sk{विपर्ययः} \sk{न} \sk{अभविष्यत्}, \sk{तर्हि} \sk{अहं} \sk{समये} \sk{आगमिष्यम्।}
\end{quote}

\begin{quote}
\emph{yadi diśāyāḥ viparyayaḥ na abhaviṣyat, tarhi ahaṃ samaye āgamiṣyam.}
\end{quote}

\begin{quote}
If I had not mixed up the direction, I would have come on time.
\end{quote}

A masculine noun.

\begin{cousinweb}
\textbf{\sk{विपर्ययः}} (\emph{viparyayaḥ}) is built on \emph{vi}, \emph{pari} and the root √i, \textquotedblleft{}to go\textquotedblright{}: a going round the wrong way; MW gives \emph{viparyaya} as reversed, inverted, contrary (MW).
\end{cousinweb}

\begin{grammarlens}[title={What you've built}]
One more word for saying what happens if.
\end{grammarlens}

\subsection*{Your turn}
\begin{itemize}
\raggedright
  \item \emph{Say it:} \emph{viparyayaḥ}
  \item \emph{Say it:} \emph{viparyayaḥ}, once more
  \item \emph{Say it:} \emph{yadi diśāyāḥ viparyayaḥ na abhaviṣyat, tarhi ahaṃ samaye āgamiṣyam.} — If I had not mixed up the direction, I would have come on time.
  \item \emph{Recall:} say the Sanskrit for come now, then the Sanskrit for guards well, then say \emph{viparyayaḥ} again
\end{itemize}

\begin{tcolorbox}[breakable,colback=teal!4,colframe=teal!35!black,title={Before you move on}]
What does \sk{विपर्ययः} mean? (\textquotedblleft{}A reversal; a mix-up\textquotedblright{}.) Say it once more.
\end{tcolorbox}
\section[\texorpdfstring{\sk{उपालम्भः} (upālambhaḥ)}{उपालम्भः (upālambhaḥ)}]{\sk{उपालम्भः} (upālambhaḥ) — a reproach}

\label{lesson:SA-C553-upalambhah}



\begin{quote}
Before the new one: say the Sanskrit for guards well, then the Sanskrit for a reversal.
\end{quote}

\subsection*{You'll want to know: \sk{उपालम्भः}}
\textbf{\sk{उपालम्भः}} — \emph{upālambhaḥ} — \textquotedblleft{}a reproach\textquotedblright{}.

\begin{quote}
\sk{यदि} \sk{अहं} \sk{सत्यम्} \sk{अवदिष्यम्}, \sk{तर्हि} \sk{उपालम्भः} \sk{न} \sk{अभविष्यत्।}
\end{quote}

\begin{quote}
\emph{yadi ahaṃ satyam avadiṣyam, tarhi upālambhaḥ na abhaviṣyat.}
\end{quote}

\begin{quote}
If I had told the truth, there would have been no reproach.
\end{quote}

A masculine noun.

\begin{cousinweb}
\textbf{\sk{उपालम्भः}} (\emph{upālambhaḥ}) is built on \emph{upa}, \emph{ā} and the root √labh, \textquotedblleft{}to take hold of\textquotedblright{}, which also gives \emph{labhate}; MW lists \emph{upālambha} under \emph{upālabh}, \textquotedblleft{}to censure, to reproach\textquotedblright{}, and gives it as reproach, censure, finding fault (MW).
\end{cousinweb}

\begin{grammarlens}[title={What you've built}]
One more word for saying what happens if.
\end{grammarlens}

\subsection*{Your turn}
\begin{itemize}
\raggedright
  \item \emph{Say it:} \emph{upālambhaḥ}
  \item \emph{Say it:} \emph{upālambhaḥ}, once more
  \item \emph{Say it:} \emph{yadi ahaṃ satyam avadiṣyam, tarhi upālambhaḥ na abhaviṣyat.} — If I had told the truth, there would have been no reproach.
  \item \emph{Recall:} say the Sanskrit for guards well, then the Sanskrit for a reversal, then say \emph{upālambhaḥ} again
\end{itemize}

\begin{tcolorbox}[breakable,colback=teal!4,colframe=teal!35!black,title={Before you move on}]
What does \sk{उपालम्भः} mean? (\textquotedblleft{}A reproach\textquotedblright{}.) Say it once more.
\end{tcolorbox}
